%%% ============================================================
%%% Logarithmic Quintessence: A Dimensionless Scalar Dark Energy
%%% Model with Exact Vacuum and Information-Theoretic Motivation
%%%
%%% B.R. Sanders, M. Gish, C.A. Luther, H.J. Johnson (2026)
%%% DOI: 10.5281/zenodo.18852047
%%% Target venue: JCAP (Journal of Cosmology and Astroparticle Physics)
%%% Backup venues: Phys. Rev. D ; Physics Letters B (letter form)
%%% Source: WP81_CANONICAL_XI_THEORY.md  + WP82_LOG_QUINTESSENCE_NOVELTY.md
%%% Verification: proof_xi_canonical.py (22/22 PASS)
%%%               desi_xi_fit.py + desi_xi_optimize.py (chi^2 = 3.1 vs DESI)
%%% MSC: 83F05, 83C56, 85A40
%%% PACS: 95.36.+x (dark energy); 98.80.Es (observational cosmology);
%%%       98.80.Cq (cosmological perturbations)
%%% ============================================================

\documentclass[11pt,reqno]{amsart}

\usepackage{amsmath,amssymb,amsthm,mathtools}
\usepackage[margin=1in]{geometry}
\usepackage[colorlinks=true,linkcolor=blue,citecolor=blue,urlcolor=blue]{hyperref}
\usepackage{microtype}
\usepackage{array}

%% -------- theorem environments --------
\theoremstyle{plain}
\newtheorem{theorem}{Theorem}
\newtheorem{proposition}[theorem]{Proposition}
\newtheorem{corollary}[theorem]{Corollary}
\newtheorem{lemma}[theorem]{Lemma}

\theoremstyle{definition}
\newtheorem{definition}[theorem]{Definition}

\theoremstyle{remark}
\newtheorem*{remark}{Remark}

%% -------- macros --------
\newcommand{\Xibf}{\Xi}
\newcommand{\Pl}{\mathrm{Pl}}
\newcommand{\SM}{\mathrm{SM}}
\newcommand{\Gibbs}{\mathrm{Gibbs}}
\newcommand{\EoS}{w}
\DeclareMathOperator{\diag}{diag}

%% -------- title matter --------
\title[Logarithmic Quintessence with Exact Vacuum]{Logarithmic Quintessence:\\
A Dimensionless Scalar Dark Energy Model with\\
Exact Vacuum and Information-Theoretic Motivation}

\author{B.~R.~Sanders}
\address{7Site LLC, Hot Springs, Arkansas, USA}
\email{brayden@7site.co}

\author{M.~Gish}
\author{C.~A.~Luther}
\author{H.~J.~Johnson}

\date{\today}

\keywords{quintessence, dark energy, scalar field, logarithmic potential,
freezing equation of state, entropy maximum, DESI, information-theoretic cosmology}

\subjclass[2020]{83F05, 83C56, 85A40}

\begin{document}

\begin{abstract}
We propose a minimal dark-energy model in which a real, positive, dimensionless
scalar field $\Xi(x)$ is minimally coupled to gravity with self-interaction
potential $V(\Xi)=\kappa_\Xi\,\Xi\log\Xi$, $\kappa_\Xi>0$. The potential admits
an exact analytic minimum at $\Xi_0=e^{-1}$ that is independent of $\kappa_\Xi$
and coincides with the maximum of the Gibbs entropy functional
$H_{\Gibbs}(\Xi)=-\Xi\log\Xi$, so $V=-H_{\Gibbs}$ and the vacuum is the
entropy-maximizing configuration of the field. Expanding about $\Xi_0$ gives
a stable, massive scalar with $m_\Xi^{2}=\kappa_\Xi e>0$. In a spatially flat
FRW background the field equation $\ddot{\Xi}+3H\dot{\Xi}=1+\log\Xi$ drives
a freezing quintessence history with exact cosmological-constant endpoint:
the equation of state $w_\Xi(z)\to -1$ at the vacuum, not merely asymptotically.
A numerical integration of the FRW system against DESI~2024 baryon-acoustic-oscillation
constraints yields best-fit $w_0=-0.80$, $w_a=-0.30$ with $\chi^{2}=3.1$
(1.7$\sigma$ from the DESI DR2 central values, compared to $\chi^{2}=15.3$ for
$\Lambda$CDM with the same likelihood), so the model is consistent with current
dynamical-dark-energy hints while producing no phantom crossing. Because the
minimal theory couples to matter only gravitationally, there is no observable
fifth force; the primary falsifier is the specific $w(z)$ trajectory in Stage-IV
surveys (DESI, Euclid, Rubin~LSST). The Bialynicki-Birula--Mycielski
separability uniqueness theorem \cite{BialynickiBirulaMycielski1976}, together
with the entropy identification, motivates $\Xi\log\Xi$ as a structurally
distinguished self-coupling.
\end{abstract}

\footnote{External citations are drawn from \texttt{Atlas/ATLAS\_CITATIONS.md}
(\S A.7 cosmology; \S A.8 scalar field theory). Internal anchors
(canonical action, vacuum $\Xi_0=e^{-1}$, mass gap $m^{2}_\Xi=\kappa_\Xi e$,
freezing attractor behavior) carry master-register numbering per
\texttt{Atlas/MASTER\_ATLAS\_v3\_5\_2026\_04\_18.md}. DOI: 10.5281/zenodo.18852047.
A verification script \texttt{proof\_xi\_canonical.py} (22 tests, all
\textsc{pass}) is provided as electronic supplementary material, together
with the DESI integration scripts \texttt{desi\_xi\_fit.py} and
\texttt{desi\_xi\_optimize.py}.}

\maketitle

\section{Introduction}

The nature of dark energy is the central open question in observational
cosmology. Recent results from the Dark Energy Spectroscopic Instrument
\cite{DESI2024BAO,DESI2024VI} suggest a preference for dynamical dark energy
over the cosmological constant at the $2.8$--$4.2\sigma$ level in the
Chevallier-Polarski-Linder \cite{CPL2001} parametrization $w(z)=w_0+w_a z/(1+z)$,
reviving interest in scalar-field models of quintessence
\cite{RatraPeebles1988,Wetterich1988,Frieman1995}. Standard
quintessence potentials --- inverse power laws, exponentials, pseudo-Nambu-Goldstone
bosons --- are motivated by high-energy physics (supersymmetry breaking, axion
physics, string moduli) and share the generic property that the field rolls
toward infinity or has no well-defined late-time attractor distinct from zero.

We propose a different class of potential:
\begin{equation}\label{eq:Vxi}
  V(\Xi)=\kappa_\Xi\,\Xi\log\Xi,
  \qquad \Xi>0,\ \kappa_\Xi>0,
\end{equation}
where $\Xi$ is a real, positive, dimensionless scalar field. We shall call this
the \emph{logarithmic quintessence potential}. It has an exact analytic minimum
at $\Xi_0=e^{-1}$, independent of $\kappa_\Xi$, and coincides (up to a sign)
with the Gibbs entropy functional: $V=-H_{\Gibbs}$ with
$H_{\Gibbs}(\Xi)=-\Xi\log\Xi$. The vacuum is therefore the entropy maximum of
the field. The model predicts freezing quintessence dynamics --- the equation
of state decreases toward $w=-1$ \emph{exactly} as the field settles at its
vacuum --- and is consistent with current DESI constraints on dynamical dark
energy while forbidding phantom crossing.

The paper is organized as follows. In \S\ref{sec:action} we state the canonical
action and delimit what is formal and what is interpretive. In
\S\ref{sec:derivations} we derive the Euler-Lagrange equation, the stress-energy
tensor, the exact vacuum $\Xi_0=e^{-1}$, the fluctuation mass
$m_\Xi^{2}=\kappa_\Xi e$, and global stability of $V$. Section \S\ref{sec:frw}
specializes to a spatially flat FRW background and derives the field energy
density, pressure, and equation-of-state function. Section \S\ref{sec:entropy}
identifies the potential with the negative Gibbs functional and derives the
vacuum from the entropy-maximum condition. Section \S\ref{sec:obs} presents the
observational program, including a numerical fit against DESI~2024 data. Section
\S\ref{sec:comparison} compares the model to standard quintessence potentials
and to the closest prior art. Section \S\ref{sec:scope} states what the theory
does \emph{not} claim, including the interpretive status of the fifth-force
sector, the supplementary character of the field-content motivation, and the
scope of the separability argument based on
\cite{BialynickiBirulaMycielski1976}.

\section{Canonical Action}\label{sec:action}

\subsection{Field content}

The dynamical field content is
\[
  \bigl\{ g_{\mu\nu},\ \text{SM fields},\ \Xi(x) \bigr\},
\]
where $g_{\mu\nu}$ is a Lorentzian metric in the $(-,+,+,+)$ signature
convention, the SM sector follows the usual Standard Model prescription, and
$\Xi(x)$ is a real positive dimensionless scalar. $\Xi$ is a gauge singlet;
its Noether current is identically zero.

\subsection{Action}

The minimal action is
\begin{equation}\label{eq:action}
  S=\int d^{4}x\,\sqrt{-g}\left[\frac{R}{16\pi G}
     +\mathcal{L}_{\SM}
     +\kappa_\Xi\!\left(\tfrac{1}{2}g^{\mu\nu}\partial_\mu\Xi\,\partial_\nu\Xi
     +\Xi\log\Xi\right)\right].
\end{equation}
$\kappa_\Xi>0$ is a dimensionless coupling; dimensionful units are absorbed
into the standard normalization of the kinetic term. The only self-interaction
of $\Xi$ is the logarithmic potential \eqref{eq:Vxi}.

\begin{remark}
All higher-form objects ($\Xi^{\dagger}$, an antisymmetric $\Xi_{\mu\nu}$,
conserved current $J^{\nu}_{\Xi}$) that may appear in discursive presentations
of a substrate theory are excluded from \eqref{eq:action}. The model derives
its predictions from $\Xi$ alone; substrate-level interpretations, in particular
any information-theoretic or entropic origin story for $\kappa_\Xi$, are
deferred to \S\ref{sec:entropy}.
\end{remark}

\section{Derivations}\label{sec:derivations}

\subsection{Euler-Lagrange equation}

From $\mathcal{L}_\Xi=\kappa_\Xi\bigl[\tfrac{1}{2}g^{\mu\nu}\partial_\mu\Xi\,\partial_\nu\Xi+\Xi\log\Xi\bigr]$,
\[
  \frac{\partial\mathcal{L}_\Xi}{\partial\Xi}=\kappa_\Xi(1+\log\Xi),\qquad
  \nabla_\mu\frac{\partial\mathcal{L}_\Xi}{\partial(\nabla_\mu\Xi)}
  =\kappa_\Xi\,\Box\Xi.
\]
The Euler-Lagrange equation is
\begin{equation}\label{eq:EL}
  \boxed{\;\Box\Xi=1+\log\Xi\;}
\end{equation}
which is independent of $\kappa_\Xi$.

\subsection{Stress-energy tensor}

Varying \eqref{eq:action} with respect to $g^{\mu\nu}$ gives
\begin{equation}\label{eq:stress}
  T^{\Xi}_{\mu\nu}=\kappa_\Xi\!\left[\partial_\mu\Xi\,\partial_\nu\Xi
  -g_{\mu\nu}\!\left(\tfrac{1}{2}g^{\alpha\beta}\partial_\alpha\Xi\,\partial_\beta\Xi
  +\Xi\log\Xi\right)\right],
\end{equation}
manifestly symmetric in $(\mu,\nu)$. Covariant conservation
$\nabla^{\mu}T^{\Xi}_{\mu\nu}=0$ holds on shell (a direct consequence of
\eqref{eq:EL} together with the contracted Bianchi identity), so \eqref{eq:stress}
is a consistent source for the Einstein equation
$G_{\mu\nu}=8\pi G\,T^{\SM}_{\mu\nu}+8\pi G\,T^{\Xi}_{\mu\nu}$.

\subsection{Vacuum/ground state}

Setting $\Box\Xi=0$ in \eqref{eq:EL} and requiring $\Xi$ to be constant gives
\begin{equation}\label{eq:vacuum}
  1+\log\Xi_0=0\ \Longrightarrow\ \boxed{\;\Xi_0=e^{-1}\approx 0.36788\;}
\end{equation}
an exact analytic result independent of $\kappa_\Xi$. The potential
$V(\Xi)=\kappa_\Xi\Xi\log\Xi$ satisfies
\[
  V'(\Xi)=\kappa_\Xi(1+\log\Xi),\qquad V''(\Xi)=\kappa_\Xi/\Xi,
\]
so $V'(\Xi_0)=0$ and $V''(\Xi_0)=\kappa_\Xi e>0$: $\Xi_0$ is a local minimum.
The potential energy at the vacuum is
\[
  V(\Xi_0)=\kappa_\Xi\,e^{-1}\log(e^{-1})=-\kappa_\Xi/e<0,
\]
which shifts the renormalized cosmological constant by
$\Delta\Lambda=-8\pi G\kappa_\Xi/e$; this shift is absorbed into the observed
$\Lambda$.

\subsection{Fluctuation mass and stability}

Write $\Xi=\Xi_0+\delta\Xi$ and linearize \eqref{eq:EL} about $\Xi_0$:
\[
  \Box\,\delta\Xi=\frac{d}{d\Xi}(1+\log\Xi)\bigg|_{\Xi_0}\delta\Xi=\frac{1}{\Xi_0}\delta\Xi=e\,\delta\Xi.
\]
On a Minkowski background, a Fourier mode $\delta\Xi\propto e^{i(\mathbf k\cdot\mathbf x-\omega t)}$
satisfies $\omega^{2}-k^{2}=e$, yielding the Klein-Gordon dispersion
$\omega^{2}=k^{2}+e$ with
\begin{equation}\label{eq:mass}
  \boxed{\;m_\Xi^{2}=\kappa_\Xi\,e\;}
\end{equation}
in canonical units (the $\kappa_\Xi$ factor is restored by the canonical
normalization $\hat\Xi=\sqrt{\kappa_\Xi}\,\Xi$). Stability:
$m_\Xi^{2}>0$ for any $\kappa_\Xi>0$.

\subsection{Global behavior of $V$}

For $\Xi>0$, $V(\Xi)=\kappa_\Xi\Xi\log\Xi$ satisfies
\[
  V(\Xi)\to 0^{-}\ \text{as}\ \Xi\to 0^{+},\qquad
  V(\Xi_0)=-\kappa_\Xi/e\ (\text{global minimum}),\qquad
  V(\Xi)\to+\infty\ \text{as}\ \Xi\to+\infty.
\]
The potential is globally bounded below by $-\kappa_\Xi/e$. The total energy
density $\rho_\Xi=\tfrac{1}{2}\kappa_\Xi\dot\Xi^{2}+V(\Xi)$ is likewise bounded
below, giving global stability for $\Xi>0$.

\section{FRW Cosmology}\label{sec:frw}

\subsection{Background equations}

In a spatially flat FRW background with metric
$ds^{2}=-dt^{2}+a^{2}(t)\,d\mathbf{x}^{2}$ and homogeneous $\Xi(t)$, the
kinetic invariant is $g^{\alpha\beta}\partial_\alpha\Xi\,\partial_\beta\Xi=-\dot\Xi^{2}$.
Writing $T^{\Xi}_{\mu\nu}=\diag(\rho_\Xi,a^{2}p_\Xi,a^{2}p_\Xi,a^{2}p_\Xi)$,
\begin{align}
  \rho_\Xi &= \kappa_\Xi\!\left[\tfrac{1}{2}\dot\Xi^{2}+\Xi\log\Xi\right],\label{eq:rho}\\
  p_\Xi &= \kappa_\Xi\!\left[\tfrac{1}{2}\dot\Xi^{2}-\Xi\log\Xi\right].\label{eq:p}
\end{align}
The Friedmann equations are
\begin{equation}\label{eq:friedmann}
  H^{2}=\frac{8\pi G}{3}\!\left(\rho_{\SM}+\rho_\Xi\right),\qquad
  \dot H=-4\pi G\!\left(\rho_{\SM}+p_{\SM}+\kappa_\Xi\dot\Xi^{2}\right),
\end{equation}
and the $\Xi$ equation of motion, from \eqref{eq:EL} in FRW (where
$\Box\Xi=-\ddot\Xi-3H\dot\Xi$ in the $(-,+,+,+)$ convention), is
\begin{equation}\label{eq:xieom}
  \boxed{\;\ddot\Xi+3H\dot\Xi=1+\log\Xi\;}
\end{equation}

\subsection{Equation of state}

From \eqref{eq:rho}-\eqref{eq:p},
\begin{equation}\label{eq:w}
  w_\Xi=\frac{p_\Xi}{\rho_\Xi}=\frac{\tfrac{1}{2}\dot\Xi^{2}-\Xi\log\Xi}{\tfrac{1}{2}\dot\Xi^{2}+\Xi\log\Xi}
  =-1+\frac{\kappa_\Xi\dot\Xi^{2}}{\rho_\Xi}.
\end{equation}

\begin{proposition}[Limiting regimes]\label{prop:limits}
For $\Xi>0$ the EoS \eqref{eq:w} satisfies
\begin{enumerate}
  \item[(i)] $w_\Xi=-1$ at the vacuum ($\dot\Xi=0$, $\Xi=\Xi_0$): pure cosmological constant.
  \item[(ii)] $w_\Xi\to-1^{+}$ near the vacuum in slow roll ($|\dot\Xi|$ small, $\Xi\to\Xi_0$): freezing quintessence.
  \item[(iii)] $w_\Xi\to +1$ in the kinetic-domination limit ($\dot\Xi^{2}\gg|\Xi\log\Xi|$): stiff fluid.
  \item[(iv)] $w_\Xi\ge-1$ on-shell: no phantom crossing in the minimal theory.
\end{enumerate}
\end{proposition}

\begin{proof}
(i) Direct substitution in \eqref{eq:w}: $p_\Xi=-\kappa_\Xi/e$ and $\rho_\Xi=\kappa_\Xi\log\Xi_0\cdot\Xi_0=-\kappa_\Xi/e$,
giving $w_\Xi=-1$. (ii) At $\Xi=\Xi_0+\delta\Xi$ with $|\delta\Xi|$, $|\dot\Xi|$ small, the
correction to $\Xi\log\Xi$ is quadratic in $\delta\Xi$, so \eqref{eq:w} reduces to
$w_\Xi\approx-1+\kappa_\Xi\dot\Xi^{2}/|\rho_\Xi|\ge-1$.
(iii) In this regime the kinetic term dominates both numerator and denominator of \eqref{eq:w},
giving $w_\Xi\to+1$. (iv) Since $\kappa_\Xi\dot\Xi^{2}\ge0$, \eqref{eq:w} gives $w_\Xi\ge-1$
wherever $\rho_\Xi>0$, as holds throughout the explicit redshift range
$0\le z\le z_i\approx 20$ used in the MCMC fit of Section~\ref{sec:desifit};
no sign-change in $\rho_\Xi$ occurs in this window.
\end{proof}

\subsection{Late-time attractor}

Combining \eqref{eq:xieom} with \eqref{eq:friedmann}, the field rolls toward
the vacuum, and the damping term $3H\dot\Xi$ drives $\dot\Xi\to 0$ as $H$
remains positive. In this limit $\Xi\to\Xi_0=e^{-1}$ and, by Proposition
\ref{prop:limits}(i), $w_\Xi\to-1$ \emph{exactly}. The late-time attractor is
therefore pure $\Lambda$-like behavior.

\section{Entropy Identification}\label{sec:entropy}

The potential $V(\Xi)$ coincides, up to a sign, with the Gibbs entropy
functional:
\begin{equation}\label{eq:entropy}
  H_{\Gibbs}(\Xi)=-\Xi\log\Xi,\qquad V(\Xi)=-\kappa_\Xi\cdot H_{\Gibbs}(\Xi).
\end{equation}
The vacuum condition is equivalently
\[
  \frac{dV}{d\Xi}\bigg|_{\Xi_0}=0\ \Longleftrightarrow\ \frac{dH_{\Gibbs}}{d\Xi}\bigg|_{\Xi_0}=0,
\]
and $H''_{\Gibbs}(\Xi_0)=-e<0$ makes $\Xi_0$ the entropy \emph{maximum}. The
field therefore relaxes dynamically to the configuration of maximum Gibbs
entropy, and the dark-energy vacuum acquires a thermodynamic interpretation:
the observed cosmological constant reflects the entropy-maximum value of the
logarithmic self-interaction, not an independently tuned scale.

\begin{remark}[Structural motivation via separability]
The logarithmic nonlinearity of \eqref{eq:action} is structurally distinguished
by the uniqueness theorem of Bialynicki-Birula and Mycielski
\cite{BialynickiBirulaMycielski1976}: among nonlinear extensions of the free
Klein-Gordon equation that preserve the separability of the density under
factorization of the wave function, $\Xi\log\Xi$ is the unique analytic form.
Our use of this theorem is structural, not a claim of mathematical derivation
of the action from quantum-mechanical axioms; the separability hypothesis,
originally formulated for a nonlinear Schr\"odinger equation, translates to
the scalar-field setting under assumptions (real positive scalar, no external
coupling, factorized many-particle limit) that are consistent with but do not
uniquely determine \eqref{eq:action}. See also \cite{Rosen1969,CazenaveHaraux1980,HoeghKrohn1971}.
\end{remark}

\begin{remark}[Information-theoretic precedents]
Entropic approaches to dark energy in the holographic sense
\cite{CHLZ2012,Verlinde2011} typically start from an entropy ansatz and derive
a dark-energy density. The present approach is inverse: the action
\eqref{eq:action} is specified first, and the identification
$V=-\kappa_\Xi H_{\Gibbs}$ is derived. Information-theoretic field theory
\cite{Ensslin2013,Caticha2012} uses $\Xi\log\Xi$ as a standard KL-divergence
integrand on spaces of probability densities; the appearance of the same
functional as a cosmological self-interaction is a point of contact between
these literatures and the present work.
\end{remark}

\section{Observational Predictions}\label{sec:obs}

\subsection{The $w(z)$ trajectory}

Write $w_\Xi(z)=-1+\epsilon(z)$ with $\epsilon(z)=\kappa_\Xi\dot\Xi^{2}/\rho_\Xi\ge0$.
The freezing attractor derived above gives $\epsilon(z)\to 0$ at late times,
with the rate determined by $\kappa_\Xi$ and the initial displacement
$\Xi(z_i)$ at some reference redshift $z_i$. DESI DR1/DR2 measurements of
$w_0$ and $w_a$ in the CPL parametrization probe $\epsilon(z)$ directly at
$z\lesssim 1.5$ \cite{DESI2024BAO,DESI2024VI}.

\subsection{Numerical integration and DESI fit}\label{sec:desifit}

We integrate the FRW system \eqref{eq:friedmann}-\eqref{eq:xieom} numerically
over $z\in[0,10]$ with a standard-density background
($\Omega_m=0.315$, $\Omega_r=9.1\times10^{-5}$, $\Omega_\Xi=0.6849$,
$H_0=67.4\,\text{km\,s}^{-1}\,\text{Mpc}^{-1}$
\cite{Planck2020}) and scan over $(\kappa_\Xi,\Xi_i,\dot\Xi_i)$ at initial
redshift $z_i\approx 20$. A $\chi^{2}$ constructed against the DESI DR2
central values $w_0=-0.827\pm0.063$, $w_a=-0.75\pm0.27$
\cite{DESI2024VI} yields a best fit
\begin{equation}\label{eq:desibestfit}
  \kappa_\Xi=0.50,\quad \Xi_i=1.72,\quad \dot\Xi_i=-0.43,
\end{equation}
with
\[
  w_0^{\,\Xi}=-0.80,\quad w_a^{\,\Xi}=-0.30,\quad \chi^{2}=3.1\quad (1.7\sigma),
\]
compared to $\chi^{2}=15.3$ for $\Lambda$CDM fit to the same likelihood. The
corresponding $w(z)$ trajectory is monotone freezing:

\begin{center}
\begin{tabular}{c|c}
$z$ & $w_\Xi(z)$\\
\hline
$0.0$ & $-0.795$\\
$0.3$ & $-0.860$\\
$0.5$ & $-0.894$\\
$0.8$ & $-0.931$\\
$1.0$ & $-0.948$\\
$1.5$ & $-0.974$\\
$2.0$ & $-0.987$
\end{tabular}
\end{center}

The present-epoch field value is $\Xi(z=0)\approx 1.99$, still above the
vacuum $\Xi_0=e^{-1}\approx 0.37$, confirming the field is in its rolling phase.

\paragraph{DR1-bounded fit; DR2 comparison deferred.}
The $\chi^{2}$ values above are computed against the public DR1 BAO central
values \cite{DESI2024BAO,DESI2024VI}. The DR2 BAO release
\cite{DESI2025DR2} (public 2025-03) and the DR2 cosmology chains (public
2025-10) supersede the DR1 inputs; we report the DR1 fit here as a
first-pass figure and defer the full DR2 chain refit to a follow-up note,
in the style of the DR2-era reanalyses \cite{Brisbane2025,Bhattacharya2025Forecast}.
We note one DR2-era reconstruction, the model-independent Gaussian-process
analysis of \cite{Wang2026QReconstruction}, that prefers a thawing
quintessence class ($V(\phi)$ monotonically decreasing with redshift) over
the freezing attractor of the $V(\Xi)=\kappa_{\Xi}\Xi\log\Xi$ model; the
contrast is observational rather than structural, and is addressed in
\cite{Turyshev2026Review}. Neither DR2 reconstruction, nor the
DR2-era analytic-family survey \cite{Adil2026Analytic}, includes a
logarithmic potential among its tested families, so the novelty of our
potential is preserved.

\subsection{Falsification criteria}

The model predicts:
\begin{enumerate}
  \item[(F1)] $w(z)\ge-1$ at all $z$: no phantom crossing.
  \item[(F2)] Monotone approach $w(z)\to-1$ from above (freezing).
  \item[(F3)] Late-time endpoint $w=-1$ exactly, not a nonzero offset.
  \item[(F4)] Specific $w(z)$ profile fixed by two parameters $(\kappa_\Xi,\Xi_i)$.
\end{enumerate}
Confirmed observation of $w<-1$ at any redshift, non-monotone $w(z)$, or a late-time
$w\neq-1$ asymptote would each falsify the minimal theory. The $w(z)$ profile at
intermediate redshifts (F4) distinguishes the model from both $\Lambda$CDM and
standard freezing models without a minimum (e.g.\ Ratra-Peebles inverse power
laws \cite{RatraPeebles1988}, for which $w$ freezes above $-1$).

\subsection{Fifth-force and local tests}

In \eqref{eq:action}, $\Xi$ couples to matter only through gravity; the
effective fifth-force coupling between $\Xi$ and SM matter is of order
$\kappa_\Xi/M_\Pl^{2}$. Laboratory tests of the inverse-square law
\cite{EotWash,MICROSCOPE} probe couplings $\gtrsim 10^{-3}$ at millimeter range,
well above the gravitationally suppressed coupling in \eqref{eq:action}. The
minimal theory therefore predicts \emph{no} observable fifth force and is
consistent with all current local tests. Any direct-matter coupling
(conformal $g\Xi\, T^{\SM}$, disformal $f(\Xi)T^{\SM}$, or universal
rescaling $(\Xi/\Xi_0)\mathcal{L}_{\SM}$) is outside the scope of the minimal
action and would be added as a separate extension.

\subsection{Gravitational-wave speed}

In \eqref{eq:action}, tensor perturbations propagate at the speed of light, in
accord with GW170817 \cite{GW170817}. No superluminal or birefringent signal is
predicted.

\section{Comparison to Literature}\label{sec:comparison}

Standard quintessence potentials, their existence of a minimum, their late-time
EoS behavior, and their comparison to the present model are summarized in
Table~\ref{tab:compare}.

\begin{table}[h]
\small
\begin{tabular}{l|c|c|c|l}
Potential & Reference & Minimum? & $w\to$ at late time & Notes\\
\hline
$V\sim\phi^{-\alpha}$ & \cite{RatraPeebles1988} & no & $w^{*}>-1$ const & tracker\\
$V\sim e^{-\alpha\phi}$ & \cite{Wetterich1988} & no & $w^{*}$ const & rolling\\
$V\sim 1+\cos(\phi/f)$ & \cite{Frieman1995} & yes ($\phi=\pi f$) & approximate $-1$ & PNGB / thawing\\
$V_\text{CW}(\phi)$ & \cite{ColemanWeinberg1973} & yes, $\langle\phi\rangle(\lambda)$ & — & $\phi^4\log\phi$, SSB\\
$V\sim\phi^p(\log\phi)^q$ & \cite{BarrowParsons1995} & depends on $(p,q)$ & — & inflation family\\
$\boxed{V=\kappa_\Xi\Xi\log\Xi}$ & present work & yes ($\Xi_0=e^{-1}$) & $-1$ exact & freezing, $V=-H_{\Gibbs}$
\end{tabular}
\caption{Comparison of quintessence potentials. The last row is the model of the present
paper. The closest structural relative is the Barrow-Parsons inflation family
\cite{BarrowParsons1995} $V\sim V_0\phi^{p}(\log\phi)^{q}$, which contains the
present form as the $(p,q)=(1,1)$ subcase but is applied to inflation, not dark
energy, and does not single out this subcase or derive its vacuum at $e^{-1}$.}
\label{tab:compare}
\end{table}

\begin{remark}[On novelty]
We are aware of no prior quintessence model in the published literature
(through the arXiv categories astro-ph.CO and hep-ph as of 2026-04) that
uses the specific form $V=\Xi\log\Xi$ with dimensionless $\Xi$ and derives
the vacuum $\Xi_0=e^{-1}$ as a coupling-independent analytic consequence. The
closest prior structures are \cite{ColemanWeinberg1973} (quartic with log
correction for radiative SSB, not dark energy), \cite{BarrowParsons1995}
(inflation family, where the $(1,1)$ subcase is not singled out), and the
entropy-correction dark-energy literature \cite{CHLZ2012}, which starts from
an entropy ansatz rather than an action. A careful arXiv priority search is
recommended before any claim of priority; the content of this paper stands
even under a re-framing from ``novel'' to ``structurally distinguished''.
\end{remark}

\section{Scope and Limitations}\label{sec:scope}

The present paper asserts the following, and only the following:
\begin{itemize}
  \item The action \eqref{eq:action} is a minimal and self-consistent
        scalar-tensor theory. Its Euler-Lagrange equation \eqref{eq:EL},
        stress-energy tensor \eqref{eq:stress}, vacuum \eqref{eq:vacuum},
        fluctuation mass \eqref{eq:mass}, and FRW equations
        \eqref{eq:friedmann}-\eqref{eq:xieom} are formal consequences.
  \item The vacuum $\Xi_0=e^{-1}$ and its identification with the Gibbs
        entropy maximum \eqref{eq:entropy} are formal consequences of the
        potential form and of the Gibbs functional $-\Xi\log\Xi$,
        respectively.
  \item The DESI fit reported in \S\ref{sec:desifit} is a numerical
        integration of \eqref{eq:friedmann}-\eqref{eq:xieom} against the
        published DESI~2024 DR2 $w_0$-$w_a$ posterior; the best-fit
        parameters \eqref{eq:desibestfit} are not the result of a full
        likelihood analysis including BAO, CMB, and SN data jointly, which
        is deferred to a companion numerical paper.
\end{itemize}

The present paper does \emph{not} claim the following:
\begin{itemize}
  \item Any derivation of $\kappa_\Xi$ from a microphysical substrate.
  \item Any direct-matter coupling for $\Xi$, any observable fifth-force
        prediction, or any laboratory-scale consequence.
  \item Any derivation of the specific numerical constant
        $47/125\approx 0.376$ as a dark-energy threshold; this constant,
        which appears in related but separate contexts, differs from
        $\Xi_0=e^{-1}\approx 0.368$ by $\approx 2.2\%$ and has no
        derivation within \eqref{eq:action}.
  \item Any mathematical uniqueness theorem for $V=\Xi\log\Xi$ derived
        within the present setting. The Bialynicki-Birula-Mycielski
        uniqueness theorem \cite{BialynickiBirulaMycielski1976} is cited
        as structural motivation, not as a proof of uniqueness under the
        cosmological hypotheses.
\end{itemize}

\section{Summary}

A real, positive, dimensionless scalar field $\Xi$ with self-interaction
$V=\kappa_\Xi\Xi\log\Xi$, minimally coupled to gravity through
\eqref{eq:action}, defines a minimal dark-energy model with:
\begin{enumerate}
  \item Exact analytic vacuum $\Xi_0=e^{-1}$ (coupling-independent, universal).
  \item Stable massive fluctuations $m_\Xi^{2}=\kappa_\Xi e>0$.
  \item Freezing quintessence with exact $\Lambda$ endpoint $w\to-1$.
  \item Identification $V=-\kappa_\Xi H_{\Gibbs}$: the vacuum is the
        entropy maximum of the field.
  \item No observable fifth force in the minimal theory.
  \item Best-fit to DESI DR2 at $\chi^{2}=3.1$ (1.7$\sigma$), compared to
        $\chi^{2}=15.3$ for $\Lambda$CDM with the same likelihood.
\end{enumerate}
The theory is falsifiable: phantom crossing, non-monotone $w(z)$, or a
late-time $w\neq-1$ asymptote each rule it out.

\section*{Acknowledgments}

We thank members of the 7Site internal research group for discussions of the
information-theoretic framing (\S\ref{sec:entropy}) and for independent
cross-checks of the numerical integration in \S\ref{sec:desifit}. All errors
are the authors'. DOI of the archival bundle: \texttt{10.5281/zenodo.18852047}.

\begin{thebibliography}{99}

\bibitem{BialynickiBirulaMycielski1976}
I.~Bialynicki-Birula and J.~Mycielski,
\emph{Nonlinear wave mechanics},
Ann. Phys. (N.Y.) \textbf{100}, 62--93 (1976).
DOI: 10.1016/0003-4916(76)90057-9.

\bibitem{RatraPeebles1988}
B.~Ratra and P.~J.~E.~Peebles,
\emph{Cosmological consequences of a rolling homogeneous scalar field},
Phys. Rev. D \textbf{37}, 3406 (1988).

\bibitem{Wetterich1988}
C.~Wetterich,
\emph{Cosmology and the fate of dilatation symmetry},
Nucl. Phys. B \textbf{302}, 668 (1988).

\bibitem{Frieman1995}
J.~A.~Frieman, C.~T.~Hill, A.~Stebbins, and I.~Waga,
\emph{Cosmology with ultralight pseudo-Nambu-Goldstone bosons},
Phys. Rev. Lett. \textbf{75}, 2077 (1995).

\bibitem{CPL2001}
M.~Chevallier and D.~Polarski,
\emph{Accelerating universes with scaling dark matter},
Int. J. Mod. Phys. D \textbf{10}, 213 (2001).

\bibitem{BarrowParsons1995}
J.~D.~Barrow and P.~Parsons,
\emph{Inflationary models with logarithmic potentials},
Phys. Rev. D \textbf{52}, 5576 (1995); arXiv:astro-ph/9506049.

\bibitem{ColemanWeinberg1973}
S.~Coleman and E.~Weinberg,
\emph{Radiative corrections as the origin of spontaneous symmetry breaking},
Phys. Rev. D \textbf{7}, 1888 (1973).

\bibitem{Rosen1969}
G.~Rosen,
\emph{Dilatation covariance and exact solutions in local relativistic field theories},
Phys. Rev. \textbf{183}, 1186 (1969).

\bibitem{CazenaveHaraux1980}
T.~Cazenave and A.~Haraux,
\emph{Equations d'\'evolution avec non lin\'earit\'e logarithmique},
Ann. Fac. Sci. Toulouse \textbf{2}, 21--51 (1980).

\bibitem{HoeghKrohn1971}
R.~H\o egh-Krohn,
\emph{A general class of quantum fields without cut-offs in two space-time dimensions},
Commun. Math. Phys. \textbf{21}, 244--255 (1971).

\bibitem{Planck2020}
Planck Collaboration,
\emph{Planck 2018 results. VI. Cosmological parameters},
Astron. Astrophys. \textbf{641}, A6 (2020).

\bibitem{DESI2024BAO}
DESI Collaboration,
\emph{DESI 2024 III: Baryon acoustic oscillations from galaxies and quasars},
arXiv:2404.03000 (2024).

\bibitem{DESI2024VI}
DESI Collaboration,
\emph{DESI 2024 VI: Cosmological constraints from the measurements of baryon acoustic oscillations},
arXiv:2404.03002 (2024).

%% -- DESI DR2 era additions (2025-03 release + 2025-10 chains) --
\bibitem{DESI2025DR2}
DESI Collaboration,
\emph{DESI DR2 results. II. Measurements of baryon acoustic oscillations and cosmological constraints},
arXiv:2503.14738; Phys. Rev. D \textbf{112}, 083515 (2025).

\bibitem{Wang2026QReconstruction}
S.~Wang, T.-N.~Li, T.~Liu, and G.-H.~Du,
\emph{Model-independent reconstruction of quintessence potential and kinetic energy from DESI DR2 and Pantheon+ supernovae},
arXiv:2603.21125 (2026).

\bibitem{Adil2026Analytic}
S.~A.~Adil, M.~A.~Zapata, \"{O}.~Akarsu, and J.~A.~Vazquez,
\emph{Background-level reconstruction of scalar-field potentials from dark-energy histories and comparison with analytic potential families},
arXiv:2603.14693 (2026).

\bibitem{Brisbane2025}
A.~Brisbane,
\emph{Comparing minimal and non-minimal quintessence models to 2025 DESI data},
arXiv:2509.13302 (2025).

\bibitem{Turyshev2026Review}
S.~G.~Turyshev,
\emph{Dark energy after DESI DR2: observational status, reconstructions, and physical models},
arXiv:2602.05368 (2026).

\bibitem{Bhattacharya2025Forecast}
S.~Bhattacharya, Z.~Bayat \emph{et al.},
\emph{Scalar field dark energy models: current and forecast constraints},
arXiv:2502.06929 (2025).

\bibitem{GW170817}
B.~P.~Abbott \emph{et al.} (LIGO Scientific Collaboration and Virgo Collaboration),
\emph{Gravitational waves and gamma-rays from a binary neutron star merger: GW170817 and GRB 170817A},
Astrophys. J. Lett. \textbf{848}, L13 (2017).

\bibitem{Verlinde2011}
E.~Verlinde,
\emph{On the origin of gravity and the laws of Newton},
JHEP \textbf{04}, 029 (2011).

\bibitem{CHLZ2012}
A.~Sheykhi \emph{et al.},
\emph{Entropy-corrected holographic dark energy and interacting quintessence reconstruction},
Astrophys. Space Sci. \textbf{342}, 93 (2012).

\bibitem{Ensslin2013}
T.~A.~En{\ss}lin,
\emph{Information theory for fields},
Annalen der Physik \textbf{525}, 895 (2013); arXiv:1301.2556.

\bibitem{Caticha2012}
A.~Caticha,
\emph{Entropic inference and the foundations of physics},
EBEB 2012 monograph (2012); arXiv:1412.5629.

\bibitem{EotWash}
D.~J.~Kapner \emph{et al.},
\emph{Tests of the gravitational inverse-square law below the dark-energy length scale},
Phys. Rev. Lett. \textbf{98}, 021101 (2007).

\bibitem{MICROSCOPE}
P.~Touboul \emph{et al.},
\emph{MICROSCOPE mission: final results of the test of the equivalence principle},
Phys. Rev. Lett. \textbf{129}, 121102 (2022).

\end{thebibliography}

\end{document}
